% French PDF pp. 67--74; Proposition 6.3.1 continues from en-11.
(v) Every finite surjective family consisting of flat
quasi-compact morphisms is covering for (fpqc).
In particular, every faithfully flat quasi-compact morphism
is covering for (fpqc).
\end{proposition}
\begin{proof}
(i) follows from 6.2.1, (ii) from 6.2.3, taking account
of the fact that a sheaf for the Zariski topology transforms
direct sums into products.
Since every representable functor is a sheaf for (Zar)
and satisfies condition (b) of (ii) by SGA 1, VIII 5.3,
\(\mathcal T_1\) is less fine than the canonical topology,
which proves (iii).

Let us prove (iv). Let \(\{S_i\to S\}\) be a family of
morphisms as in the statement. Considering a covering of
\(S\) by affine open sets, one is immediately reduced
to the case where \(S\) is affine.\nde{The rest has been
simplified, taking advantage of the fact that one now
assumes \(S\) affine.}
First let us treat the case where the morphisms \(S_i\to S\)
are flat and open (resp. étale).
Let \(S_{ij}\) be a covering of \(S_i\) by affine open sets.
Since the morphisms in question are open, the images
\(T_{ij}\) of the \(S_{ij}\) form an open covering of \(S\).
Since \(S\) is affine, hence quasi-compact, it is covered
by a finite number of open sets \(T_{ij}\), with \((i,j)\)
ranging over a finite set \(F\).
Then \(S'=\coprod_F S_{ij}\) is affine, and the morphism
\(S'\to S\) belongs to \(P'_1\), resp. \(P'_3\), hence is
covering. Since it factors through the given family,
the latter is covering.

In the case (étf), each \(S_i\) is finite over \(S\),
hence affine; in the preceding argument, one may then
take the covering \(\{S_i\}\) of \(S_i\), and one obtains
that \(S'\to S\) belongs to \(P'_4\).

Now consider the case where the morphisms \(f_i:S_i\to S\)
are flat and locally of finite presentation.
For every \(s\in S\), there exists, by (the proof of)
EGA IV$_4$, 17.16.2, an affine subscheme \(X(s)\) of
some \(S_i\), such that \(s\in f_i(X(s))\) and the
morphism \(g_i:X(s)\to S\), the restriction of \(f_i\),
is flat, of finite presentation, and quasi-finite.
Then \(g_i(X(s))\) is an open neighborhood \(U(s)\)
of \(s\) (EGA IV$_2$, 2.4.6), and, since \(S\) is
affine, it is covered by a finite number of such open
sets \(U(s_j)\), \(j=1,\ldots,n\).
Consequently \(X'=\coprod_j X(s_j)\) is affine, and
the morphism \(X'\to S\) is surjective, flat, of
finite presentation and quasi-finite, hence belongs
to \(P'_2\).\nde{This shows that, if one denotes by
\(P''_2\) the set of finite surjective families of
flat morphisms of \(\mathcal C'\) of finite presentation,
the topology generated by \(P,P''_2\) equals
\(\mathcal T_2\). On the other hand, with the notation
at the beginning of the proof of (iv), if one takes
a covering of \(S_i\) by affine open sets of finite
presentation over \(S\), one obtains that \(S'\to S\)
belongs to \(P''_2\).}
% typo: French "présentation fini" -> "présentation finie".
This completes the proof of (iv).

Let us prove (v). Let \(\{S_i\to S\}\) be a finite
faithfully flat quasi-compact family.
Let \(T_j\) be a covering of \(S\) by affine open sets.
The \(S_{ij}=T_j\times_S S_i\) are quasi-compact,
and therefore possess finite affine open coverings
\(T_{ijk}\). Each morphism \(T_{ijk}\to T_j\) is
flat, and the family \(\{T_{ijk}\to T_j\}\) is
finite and surjective, hence covering for \(\mathcal T_1\).
The family \(\{T_{ijk}\to S\}\) is therefore also
covering, by composition.
It factors through the given family, which is
therefore also covering:
\[
\xymatrix{
S_i\ar[d] & S_{ij}\ar[l]\ar[d] & T_{ijk}\ar[l]\ar[dl]\\
S & T_j\ar[l] & }.
\]
\end{proof}

\begin{corollary}{6.3.2}\label{IV.6.3.2}
Let \((\mathrm M_i)\) be the following family of morphisms:

\((\mathrm M_1)\): faithfully flat quasi-compact morphisms.

\((\mathrm M_2)\): faithfully flat morphisms locally of finite presentation.

\((\mathrm M_3)\): surjective étale morphisms.

\((\mathrm M_4)\): surjective finite étale morphisms.
\nde{The original has been corrected by adding for
\((\mathrm M_3),(\mathrm M_4)\) the surjectivity hypothesis,
which is automatically satisfied in the other cases.}
The family \((\mathrm M_i)\) satisfies axioms (a),
(b), (c), \((\mathrm d_{\mathcal T_i})\), and
\((\mathrm e_{\mathcal T_i})\) of 4.6.3.
\end{corollary}
\begin{proof}
Indeed, for (a), (b), (c), this is classical
(EGA and SGA 1, passim).\nde{Cf. EGA I, 6.6.4
for \og quasi-compact\fg, EGA II, 6.1.5 for
\og finite\fg, EGA IV$_1$, 1.6.2 for
\og locally of finite presentation\fg, EGA IV$_2$,
2.2.13 for \og faithfully flat\fg, and
EGA IV$_4$, 17.3.3 for \og étale\fg.}
By 6.3.1 (iv) and (v), \((\mathrm M_i)\) satisfies
\((\mathrm d_{\mathcal T_i})\).
It remains to see that \((\mathrm M_i)\)
satisfies \((\mathrm e_{\mathcal T_i})\);
for this it suffices to see that \((\mathrm M_i)\)
satisfies \((\mathrm e_{\mathcal T_1})\), which
implies the others. This follows from SGA 1,
VIII (nos.\ 4 and 5).
\end{proof}
\begin{corollary}{6.3.3}\label{IV.6.3.3}
If \(X\) is a scheme and \(R\) an equivalence
relation in \(X\) of type \((\mathrm M_i)\),
\(R\) is \((\mathrm M_i)\)-effective if and
only if the quotient sheaf of \(X\) by \(R\)
for \(\mathcal T_i\) is representable, and
in this case it is represented by the quotient \(X/R\).
\end{corollary}
\begin{proof}
Indeed, this is 4.6.5.
\end{proof}

\sgasection{6.4}{Effectiveness conditions}
We now seek families (N) of morphisms
satisfying axiom \((\mathrm f_{\mathcal T})\)
of 4.7. First remark that
\((\mathrm f_{\mathcal T_1})\) implies
\((\mathrm f_{\mathcal T_i})\), so that
we may restrict ourselves to the case
of the topology (fpqc).
\begin{lemma}{6.4.1}\label{IV.6.4.1}
The following families of morphisms satisfy
axiom \((\mathrm f_{\mathcal T_1})\) of 4.7,
that is, \og descend under (fpqc)\fg:

(N): open immersions.

(N'): closed immersions.

(N''): quasi-compact immersions.
\end{lemma}
\begin{proof}
By 6.3.1 (ii), it suffices to verify that
the given families descend under the
Zariski topology and under a faithfully
flat quasi-compact morphism.
The first assertion is clear; let us
verify the second. For (N), this is
SGA 1, VIII 4.4; for (N'), it is
\loccit, 1.9. For (N''), one argues
as in \loccit, 5.5, with the aid of
the two preceding results.
\end{proof}
\begin{corollary}{6.4.2}\label{IV.6.4.2}
The same result is valid for quasi-compact open immersions.
\end{corollary}
These results allow one to apply to
the present situation the general
results of 4.7.1, 4.7.2, 5.1.8,
5.3.1, etc. Let us state one as an
example, the first.
\begin{corollary}{6.4.3}\label{IV.6.4.3}
\textup{(= 4.7.1 + 4.6.10).}
Let \(X\) be a scheme and \(R\)
an equivalence relation in \(X\).
Suppose that \(R\to X\) is faithfully
flat and quasi-compact, and that
\(R\to X\times X\) is a closed
immersion (resp. open, resp. quasi-compact,
resp. quasi-compact open).
Then the quotient sheaf \(X/R\)
is the same for the topology (fpqc)
and for the canonical topology,
and for each scheme \(S\), one has
\[
(X/R)(S)=\left\{\begin{array}{l}
\text{closed (resp. open, resp. retrocompact,}\\
\text{resp. retrocompact open) subschemes }Z
\text{ of }X_S\text{ stable under }R\times S,\\
\text{such that }Z\to S\text{ is faithfully flat quasi-compact}\\
\text{and the diagram }R_Z\rightrightarrows Z\to S\text{ is exact}
\end{array}\right\}.
\]
\nde{Recall that a subscheme \(Z\)
of a scheme \(T\) is said to be
retrocompact if the immersion
\(Z\hookrightarrow T\) is
quasi-compact, cf. EGA \(0_{\mathrm{III}}\), 9.1.1.}
\end{corollary}

\sgasection{6.5}{Principal homogeneous bundles}
Let us simply point out the terminology:
\begin{center}
\begin{tabular}{ll}
Topology & Principal homogeneous bundles\\
(fpqc) & (without further qualification)\\
(ét) & quasi-isotrivial\\
(étf) & locally isotrivial\\
(Zar) & locally trivial.
\end{tabular}
\end{center}
\sgasection{6.6}{Other topologies}
One sometimes uses other topologies
on the category of schemes.
Let us point out one: the
\emph{global finite étale topology}
(étfg), generated by the pretopology
whose covering families are
surjective families consisting of
finite étale morphisms.
It is not finer than the Zariski
topology. The corresponding
principal homogeneous bundles
are called \og isotrivial\fg.
\[
\xymatrix@R=1.3em@C=1.2em{
 & (\text{canonical})\ar@{-}[d] & \\
 & (\mathrm{fpqc})\ar@{-}[d] & \\
 & (\mathrm{fppf})\ar@{-}[d] & \\
 & (\mathrm{\acute et})\ar@{-}[d] & \\
 & (\mathrm{\acute etf})\ar@{-}[dl]\ar@{-}[dr] & \\
(\mathrm{Zar})\ar@{-}[dr] & & (\mathrm{\acute etfg})\ar@{-}[dl]\\
 & (\text{chaotic}) & }.
\]
\nde{Recall (cf. 4.4.2) that the
chaotic topology is the least fine
topology, defined by \(J(S)=\{S\}\)
for every \(S\in\Ob\mathcal C\).}

\sgasection{6.7}{Homogeneous spaces}
\nde{The following numbers have been added.}
Let \(G\) be an \(S\)-group scheme,
\(X\) an \(S\)-scheme with a group
of operators \(G\) (on the left), and
\[
\Phi:G\times_S X\longrightarrow X\times_S X
\]
the morphism of \(S\)-schemes
defined set-theoretically by
\((g,x)\mapsto(gx,x)\).
Recall (cf. 5.1.0 and III.0.1)
that \(X\) is said to be a
formally principal homogeneous
space under \(G\) if the following
equivalent conditions are satisfied:
\begin{enumeratei}
\item For every \(T\to S\), the
set \(X(T)\) is empty or principal
homogeneous under \(G(T)\).
\item \(\Phi\) is an isomorphism of \(S\)-functors.
\item \(\Phi\) is an isomorphism of \(S\)-schemes.
\end{enumeratei}
(The equivalence (i) \(\Leftrightarrow\)
(ii) is clear, and one has
(ii) \(\Leftrightarrow\) (iii),
since \(\mathcal C=\Sch_{/S}\)
is a full subcategory of
\(\widehat{\mathcal C}\).)

The definition of a formally
homogeneous space (not necessarily
principal homogeneous) is obtained
by requiring that \(\Phi\) be an
epimorphism in the category of
sheaves for an appropriate topology
\(\mathcal T\).
Indeed, the condition that \(\Phi\)
be an epimorphism of \(S\)-functors
is equivalent to saying that for
every \(T\to S\), the set \(X(T)\)
is empty or homogeneous (not
necessarily principal homogeneous)
under \(G(T)\), but this condition
is too restrictive, as the following
simple example shows.
Let \(S=\Spec\RR\), \(G=\mathbf G_{\mathrm m,\RR}\),
and \(X=\mathbf G_{\mathrm m,\RR}\),
on which \(G\) acts by \(t\cdot x=t^2x\).
Then the morphism \(\Phi\) is étale,
finite, and surjective, hence an
epimorphism in the category of
sheaves for the topology (étf)
(a fortiori, an epimorphism of
\(S\)-schemes); on the other hand,
the points \(1,-1\) of \(X(\RR)\)
are not conjugate by an element
of \(G(\RR)\), so that the
morphism \(G(\RR)\times X(\RR)\to
X(\RR)\times X(\RR)\) is not surjective.
\nde{Evidently, this difficulty
comes from the fact that if
\(\mathcal C'\) is a full
subcategory of \(\widehat{\mathcal C}\)
containing \(\mathcal C\),
for example the category
\(\widetilde{\mathcal C}_{\mathcal T}\)
of sheaves on \(\mathcal C\)
for a topology \(\mathcal T\)
less fine than the canonical
topology, and if \(f:X\to Y\)
is a morphism in \(\mathcal C\),
then the implications:
\[
\begin{gathered}
f\text{ epimorphism of }\widehat{\mathcal C}
\Rightarrow f\text{ epimorphism of }\mathcal C'\\
\Rightarrow f\text{ epimorphism of }\mathcal C
\end{gathered}
\]
are in general strict.}
One is therefore led to give the following definition:
\begin{definition}{6.7.1}\label{IV.6.7.1}
Let \(G\) be an \(S\)-group,
\(X\) an \(S\)-scheme with a
group of operators \(G\),
and \(\mathcal T\) a topology
on \(\Sch_{/S}\) less fine
than the canonical topology.
One says that \(X\) is a
\emph{formally homogeneous space}
under \(G\) (relative to the
topology \(\mathcal T\)) if
the following equivalent
conditions are satisfied:
\begin{enumeratei}
\item The morphism \(\Phi:G\times_S X\to X\times_S X\)
is an epimorphism in the category
of sheaves for the topology \(\mathcal T\).
\item For every \(T\to S\)
and \(x,y\in X(T)\), there
exists a morphism \(T'\to T\)
covering for the topology
\(\mathcal T\), and \(g\in G(T')\),
such that \(y_{T'}=g\cdot x_{T'}\).
\end{enumeratei}
\end{definition}
\begin{remark}{6.7.2}\label{IV.6.7.2}
Condition (i) implies in particular
that \(\Phi\) is a universal
effective epimorphism in \(\Sch_{/S}\)
(cf. 4.4.3).
This implies, as one easily sees,
that \(\Phi\) is surjective
(cf. 1.3, N.D.E. (3)).
\end{remark}
\begin{enonce}{Proposition and definition}{6.7.3}\label{IV.6.7.3}
\nde{Cf. [Ray70], Def. VI.1.1.}
Let \(G\) be an \(S\)-group,
\(X\) an \(S\)-scheme with a
group of operators \(G\),
and \(\mathcal T\) a topology
on \(\Sch_{/S}\) less fine
than the canonical topology.
The following conditions are equivalent:

(i) \(X\) satisfies the following two hypotheses:

(1) The morphism \(\Phi:G\times_S X\to X\times_S X\)
is covering, i.e. \(X\) is a
formally homogeneous \(G\)-space.

(2) The morphism \(X\to S\)
is also covering, that is,
locally for the topology
\(\mathcal T\), it possesses
a section (cf. RefIV.4.4.8bis.bis).
% typo?: broken printed reference retained verbatim.

(ii) \og Locally on \(S\) for
the topology \(\mathcal T\)\fg,
\(X\) is isomorphic, as a
scheme with a group of operators
\(G\), to the quotient sheaf
(for \(\mathcal T\)) of \(G\)
by a subgroup scheme \(H\),
that is, there exists a
covering family \(\{S_i\to S\}\)
such that each \(X\times_S S_i\)
represents the quotient sheaf
of \(G\times_S S_i\) by some
subgroup scheme \(H_i\).

Under these conditions, one says
that \(X\) is a \emph{homogeneous
\(G\)-space} (relative to the
topology \(\mathcal T\)).
\end{enonce}
\begin{proof}
Suppose that (ii) holds.
Set \(G_i=G\times_S S_i\),
\(X_i=X\times_S S_i\).
Then \(X_i\) possesses a section
over \(S_i\), namely the
composite of the unit section
\(\varepsilon_i:S_i\to G_i\)
and the projection
\(\pi_i:G_i\to X_i=G_i/H_i\).
Thus \(X\to S\) is covering.

On the other hand, \(\pi_i\)
is covering, hence \(\pi_i\times\pi_i\)
is too (cf. 4.2.3 (C 1)
and (C 2)), and one has
a commutative diagram:
% typo?: all Gi targets in the printed diagram retained.
\[
\xymatrix{
G_i\times_{S_i}X_i\ar[r]^{\Phi_i}
 & G_i\times_{S_i}X_i\\
G_i\times_{S_i}G_i
\ar[u]^{\mathrm{id}\times\pi_i}
\ar[r]_{\sim}^{\Psi_i}
 & G_i\times_{S_i}G_i
\ar[u]_{\pi_i\times\pi_i} }
\]
where \(\Phi_i\) is deduced
from \(\Phi\) by the base
change \(S_i\to S\), and
\(\Psi_i\) is the isomorphism
defined set-theoretically by
\((g,g')\mapsto(gg',g)\).
Then \((\pi_i\times\pi_i)\circ\Psi_i\)
is covering, hence \(\Phi_i\)
is too (4.2.3 (C 3)).
This shows that \(\Phi\) is
\og locally covering\fg, hence
covering (4.2.3 (C 5)).
This proves (ii) \(\Rightarrow\) (i).

Conversely, suppose that (i)
holds, and suppose moreover
that the structural morphism
\(X\to S\) possesses a section
\(\sigma\). By EGA I, 5.3.13,
\(\sigma\) is an immersion.
Define \(H=G\times_X S\)
by the diagram below, both
of whose squares are cartesian:
\[
\xymatrix@C=3em{
H\ar[r]\ar[d] & G\ar[r]^{\mathrm{id}_G\boxtimes\sigma}\ar[d]^\pi
 & G\times_S X\ar[d]^\Phi\\
S\ar[r]_\sigma & X\ar[r]_{\mathrm{id}_X\boxtimes\sigma}
 & X\times_S X }
\]
where \(\pi,\mathrm{id}_G\boxtimes\sigma\)
and \(\mathrm{id}_X\boxtimes\sigma\)
denote the morphisms defined
set-theoretically, for \(T\to S\)
and \(g\in G(T)\), \(x\in X(T)\), by:
\[
\begin{gathered}
\pi(g)=g\cdot\sigma_T,\qquad
(\mathrm{id}_G\boxtimes\sigma)(g)=(g,\sigma_T),\\
(\mathrm{id}_X\boxtimes\sigma)(x)=(x,\sigma_T).
\end{gathered}
\]
Then \(\pi\) is covering,
and \(H\) is a subgroup
scheme of \(G\), representing
the stabilizer
\(\underline{\mathrm{Stab}}_G(\sigma)\)
of \(\sigma\) (cf. I, 2.3.3),
that is, for every \(T\to S\), one has:
\[
H(T)=\{g\in G(T)\mid g\cdot\sigma_T=\sigma_T\}.
\]
Denote by \(G/H\) the
\emph{presheaf} \(T\mapsto G(T)/H(T)\),
and by \(a(G/H)\) the associated
sheaf for the topology \(\mathcal T\).
By the preceding argument,
one obtains a commutative
diagram of morphisms of
presheaves with a group of operators \(G\):
\[
\xymatrix{
G\ar[r]^\pi\ar[d] & X\\
G/H\ar[ur]_{\overline\pi} & },
\]
where \(\overline\pi\) is a
monomorphism. Since \(\pi\)
is covering, \(\overline\pi\)
is too, and hence, by 4.3.12,
\(\overline\pi\) induces an
isomorphism \(a(G/H)\isomto X\).
One has therefore proved that:
\emph{if \(X\) is a homogeneous
\(G\)-space such that \(X\to S\)
admits a section \(\sigma\),
then \(X\) represents the
quotient sheaf \(G/H\),
where \(H=G\times_X S\)
is the stabilizer of \(\sigma\).}

In the general case, there
exists by hypothesis a covering
family \(\{S_i\to S\}\) such
that each morphism
\(X_i=X\times_S S_i\to S_i\)
possesses a section \(\sigma_i\).
Set \(G_i=G\times_S S_i\);
then the morphism
\(\Phi_i:G_i\times_{S_i}X_i\to X_i\times_{S_i}X_i\)
deduced from \(\Phi\) by the
base change \(S_i\to S\)
is still covering.
Thus, by the preceding argument,
\(X_i\simeq G_i/H_i\), where
\(H_i\) is the stabilizer in
\(G_i\) of \(\sigma_i\).
This completes the proof of
the implication (i) \(\Rightarrow\) (ii).
\end{proof}

\begin{thebibliography}{TDTE I}
\bibitem[AS]{AS}
\emph{Analysis Situs}, by J. Giraud,
Sém. Bourbaki, Exp. 256, May 1963.
\bibitem[D]{D}
\emph{Méthode de la descente}, by J. Giraud,
Mém. Soc. Math. France, vol. 2 (1964),
pp. iii--viii + 1--150.
\bibitem[MA]{MA}
\emph{Grothendieck Topologies}, by M. Artin,
mimeographed notes, Harvard, 1962.
\bibitem[SGA 1]{SGA1}
\emph{Séminaire de Géométrie Algébrique
du Bois-Marie 1960--61, Revêtements
étales et groupe fondamental},
Lecture Notes in Maths. 224 (1971),
recomposed and annotated edition,
Documents Math. 3, Soc. Math. France, 2003.
\bibitem[SGA 4]{SGA4}
\emph{Séminaire de Géométrie Algébrique
du Bois-Marie 1963--1964, Théorie
des topos et cohomologie étale
des schémas}, vols. I, II, III,
Lecture Notes in Maths. 269, 270
(1972), 305 (1973).
\bibitem[TDTE I]{TDTEI}
\emph{Techniques de descente et
théorèmes d'existence en géométrie
algébrique I. Généralités.
Descente par morphismes fidèlement plats},
by A. Grothendieck, Sém. Bourbaki,
Exp. 190, December 1959.
\bibitem[Ray70]{Ray70}
\nde{This reference has been added.}
M. Raynaud, \emph{Faisceaux amples
sur les schémas en groupes et
les espaces homogènes}, Lect.
Notes Math. 119, Springer-Verlag, 1970.
\end{thebibliography}
